\documentclass[12 pt, dvipsnames]{article}        	%sets the font to 12 pt and says this is an article (as opposed to book or other documents)
\usepackage{amsfonts, amssymb, amsmath, amsthm}					% packages to get the fonts, symbols used in most math
%\usepackage{dcolumn}
\usepackage[english]{babel}
\usepackage{latexsym}
\DeclareRobustCommand{\ojoin}{\rule[0.10ex]{.3em}{.4pt}\llap{\rule[1.40ex]{.3em}{.4pt}}}
\newcommand{\leftouterjoin}{\mathrel{\ojoin\mkern-6.5mu\Join}}
\newcommand{\rightouterjoin}{\mathrel{\Join\mkern-6.5mu\ojoin}}
\newcommand{\fullouterjoin}{\mathrel{\ojoin\mkern-6.5mu\Join\mkern-6.5mu\ojoin}}

  
\input{_reviewing.tex}
%\usepackage{setspace}               		% Together with \doublespacing below allows for doublespacing of the document

\oddsidemargin=-0.5cm                 	% These three commands create the margins required for class
\setlength{\textwidth}{6.5in}         	%
\addtolength{\voffset}{-20pt}        		%
\addtolength{\headsep}{25pt}           	%



\pagestyle{myheadings}                           	% tells LaTeX to allow you to enter information in the heading
\markright{Sitt Min Oo\hfill October 7, 2024 \hfill} 	% put your name instead of Murphy Waggoner 
																									% and put the proposition number from the book
                                                	% LaTeX will put your name on the left, the date the paper 
                                                	% is generated in the middle 
                                                 	% and a page number on the right

%%%%%%%%%%% Put your definitions here
\newcounter{defn}
\newcounter{exmp}
\newcounter{thm}

\theoremstyle{definition}
\newtheorem{definition}[defn]{Definition}


\newtheorem{theorem}[thm]{Theorem}
\newtheorem{example}[exmp]{Example}

%%%%%%%%%%% End of definitions



\newcommand{\eqn}[0]{\begin{array}{rcl}}%begin an aligned equation - allows for aligning = or inequalities.  Always use with $$ $$
\newcommand{\eqnend}[0]{\end{array} }  	%end the aligned equation

\newcommand{\qed}[0]{$\square$}        	% make an unfilled square the default for ending a proof

%\doublespacing                         	% Together with the package setspace above allows for doublespacing of the document

\begin{document}												% end of preamble and beginning of text that will be printed


\section{Preliminaries}
\label{sec:preliminaries}
Given the following pairwise disjoint infinite sets: $V$ (variables), $I$
(IRIs), $B$ (RDF blank nodes), and $L$ (RDF literals).
Following the definitions in SPARQL,
a \emph{solution mapping} $\mu$ is a mapping from variables $V$ to
associated data values of type $I \cup B \cup L$.
More formally, it is defined as a partial function $\mu: V \rightarrow T$, with
$T = I \cup B \cup L$.

Two solution mappings $\mu_1$ and $\mu_2$ are \textit{compatible}, if and only if,
$\forall v \in dom(\mu_1) \cap dom(\mu_2), \mu_{1}(v) = \mu_{2}(v)$; extending
this, the union of $\mu_1$ and $\mu_2$, $\mu_1 \cup \mu_2$, is also a
solution mapping.
Given two multisets of solution mappings $\Omega_1$ and $\Omega_2$, SPARQL
algebra~\cite{Perez2009SPARQLSemantics} defined the \textit{join} ($\Join$),
the \textit{union} ($\cup$), and the \textit{difference}  ($ \setminus $)
between $\Omega_1$ and $\Omega_2$ as follows~\cite{Perez2009SPARQLSemantics}:

\begin{align}
	\label{eq:sparql_union}
	\Omega_1 \cup \Omega_2           & = \{ \mu \  |\   \mu \in \Omega_1 \text{ or }  \mu \in \Omega_2\}                                                        \\
	\label{eq:sparql_join}
	\Omega_1 \Join \Omega_2          & = \{ \mu_1 \cup \mu_2 \ | \  \mu_1 \in \Omega_1, \mu_2 \in \Omega_2 \text{ and } \mu_1, \mu_2 \text{ are compatible. }\} \\
	\label{eq:sparql_diff}
	\Omega_1 \setminus \Omega_2      & = \{ \mu_1 \in \Omega_1 \ | \ \forall \mu_2 \in \Omega_2, \mu_1 \text{ and } \mu_2 \text{ are not compatible. }\}       
\end{align}



When writing sequences of the join and the union operators, parentheses can be
avoided since both operators are associative and
commutative~\cite{Perez2009SPARQLSemantics}:

\begin{align}
	\label{eq:asso_comm}
	\Omega_1 \Join \Omega_2 \cup \Omega_3 = \Omega_1 \Join (\Omega_2 \cup \Omega_3)
\end{align}



\section{Definitions}%
\label{sec:definitions}
\begin{definition}
	Given a solution mapping $\mu$.
	A \emph{mapping tuple}, $\xi$, is an extension on the solution mapping
	$\mu$ with a special variable, $v_{fragment}$, bound to a Literal value,
	with the condition that $v_{fragment} \notin V$.
	The special variable $v_{fragment}$ is used to store the information
	about the \emph{fragment} a solution mapping belongs to.
	A \emph{fragment}, in the context of this work, is an identifier used for
	grouping the mapping tuples.
    Since mapping tuple is an extension on top of the existing definition of
    solution mappings, in Section~\ref{sec:preliminaries}, the operations
    defined in Section~\ref{sec:preliminaries} also applies to the mapping
    tuple.
	\todoi{TODO: Check the following sentence again to make sure it can be called as a \emph{relation}}
	Finally, a relation denoted as $\Xi$, is defined as a relation containing
	a multi-set of mapping tuples.
	\label{def:mapping_tuple}
\end{definition}


\newpage

\todoi{TODO: Rework the following definitions of iterators according to the
discussion}
\begin{definition}
	An \textbf{iterator}, $I$, consists of one or more \textbf{fields},
	$\phi^{I}$, an iterator \emph{path}, and a \emph{reference formulation}.
	An iterator extracts a list of records from the data source according to the
	given iterator \emph{path}~\cite{delva2021RMLFields}, while the
	\emph{reference formulation} determines the data format~\cite{dimou_ldow_2014}
	of the data source.
	This is the entry point to further extract nested data structures with \emph{fields}.
	\label{def:iterator}
\end{definition}






\bibliographystyle{ios1}           % Style BST file.
\bibliography{bibliography}        % Bibliography file (usually '*.bib')


\end{document}
